%==============================================================================
% COVER LETTER FOR JOURNAL OF PHYSICAL CHEMISTRY LETTERS
%==============================================================================

\documentclass[11pt,a4paper]{letter}
\usepackage[utf8]{inputenc}
\usepackage[T1]{fontenc}
\usepackage{amsmath,siunitx}

\address{
Department of Physics\\
Faculty of Science\\
University of Yaoundé I\\
Yaoundé, Cameroon
}

\signature{Steve Cabrel Teguia Kouam}

\begin{document}

\begin{letter}{
Dr. Anne B. McCoy\\
Editor-in-Chief\\
The Journal of Physical Chemistry Letters\\
American Chemical Society
}

\opening{Dear Dr. McCoy,}

We submit ``Quantum-Coherent Spectral Engineering: Non-Markovian Dynamics in Photosynthetic Energy Transfer under Engineered Photon Baths'' for consideration as a Letter in The Journal of Physical Chemistry Letters.

\textbf{What we found:} Spectral engineering of the incident photon bath extends excitonic coherence lifetimes by \SIrange{20}{50}{\percent} and enhances energy transfer efficiency by \SI{25}{\percent} in the Fenna--Matthews--Olson complex at room temperature. Dual-band filtering at \SIlist{750;820}{\nano\meter} targets vibronic resonances, producing polaron-dressed states with reduced pure dephasing.

\textbf{Why it matters:} This work establishes \textit{spectral bath engineering} as a new paradigm for coherence control in open quantum systems. Unlike prior studies that treat environmental coupling as fixed, we show that spectral shaping can \textit{enhance} rather than suppress quantum effects. The mechanism generalizes to any system where vibronic resonances contribute to transport.

\textbf{Broader impact:} The spectral structure of illumination---not merely intensity---constitutes an independent control parameter for quantum dynamics. Applications include optimizing light-harvesting materials, designing quantum sensors, and controlling condensed-phase reactions. Commercial dual-band filters enable immediate experimental tests.

\textbf{Why JPCL:} Our work bridges quantum dynamics theory with immediate experimental testability via 2DES with spectrally shaped pulses. The concise Letter format matches the focused physical insight. We have structured the manuscript to meet the \num{2500}-word limit while maintaining scientific completeness.

\textbf{Key contributions:}
\begin{enumerate}
\item \textbf{Concept:} "Spectral bath engineering"---shaping excitation spectra to enhance quantum coherence
\item \textbf{Results:} \SIrange{20}{50}{\percent} coherence extension and \SI{25}{\percent} efficiency enhancement at \SI{295}{\kelvin}
\item \textbf{Mechanism:} Vibronic resonance selectively suppresses pure dephasing via polaron dressing
\item \textbf{Validation:} 12-test suite (\SI{1.8}{\percent} HEOM deviation), three-site model generalization, experimental feasibility analysis
\end{enumerate}

This computational study employs numerically exact PT-HOPS/SBD simulations. We provide quantitative 2DES predictions, commercial filter specifications, and proposed control experiments---all ready for laboratory implementation.

All authors have approved this manuscript. We declare no conflicts of interest.

\textbf{Suggested reviewers:}
\begin{itemize}
\item Gregory Scholes (Princeton) -- gscholes@princeton.edu
\item Alán Aspuru-Guzik (Toronto) -- alan@aspuru.com
\item Martin Plenio (Ulm) -- martin.plenio@uni-ulm.de
\item Jianshu Cao (MIT) -- jianshu@mit.edu
\item Graham Fleming (Berkeley) -- gfleming@berkeley.edu
\end{itemize}

\textbf{Excluded reviewers:} None requested.

This manuscript is not under consideration elsewhere. Upon acceptance, we will transfer copyright to ACS.

Thank you for your consideration.

\closing{Sincerely,}

\fromsig{
Steve Cabrel Teguia Kouam\\
Department of Physics\\
University of Douala, Cameroon\\
Email: steve.teguia@facsciences-uy1.cm
}

\end{letter}
\end{document}
